% sec-05: budget, milestones, partner site.  Sources: RFA-RM-27-001-v2 budget
% module, potential-partners/UC-San-Diego.  Figure 3 is a graphviz-type DAG.
\section{Budget, Milestones, and the Partner Site}

The request is \textbf{\$700,000 per year for five years}, \$3,500,000 total, no
cost share \cite{fundapp2}. The report asks that long-duration grants be fully
funded in year one; issued that way, year one closes the site agreement, the IRB
submission, and IND activation together rather than in sequence.

\begin{apptable}
\begin{tabularx}{\textwidth}{>{\raggedright\arraybackslash}p{1.35cm} >{\raggedright\arraybackslash}p{2.35cm} >{\raggedright\arraybackslash}X >{\raggedright\arraybackslash}p{2.5cm}}
\headrow \thead{Year} & \thead{Direct cost} & \thead{Milestone} & \thead{Go / no-go}\\
\midrule
1 & \$700,000 & Site agreement, IRB approval, IND activation, bench stop-latency verification & IND active and site executed \\
2 to 3 & \$1,400,000 & Dose level 1 cohort, 30-day safety review, escalation to RP2D through cohorts 2 and 3, interim advisory audit & No unexpected grade 4 event; MTD or RP2D declared \\
4 to 5 & \$1,400,000 & Full accrual to 18 participants, 90-day morbidity set, 24-month survival follow-up, public data and code release & Accrual at plan; package released with a DOI \\
\end{tabularx}
\end{apptable}
\tabcap{Five-year milestone and go/no-go schedule at \$700,000 per year. Each
gate is a stated, checkable condition, so the award can be stopped at a cohort
boundary rather than only at a renewal date.}

\subsection{UC San Diego Moores Cancer Center}

Moores Cancer Center is the intended partner of choice, approached at the
feasibility stage only: a nonconfidential overview, a joint feasibility meeting
with pancreatic surgery, gastrointestinal medical oncology, and trial
operations, then an investigator-initiated intake with a budget and accrual
estimate. UC San Diego is not described here as a partner, sponsor, site, or
endorser, because no written authorization exists.

\begin{appfloat}[!tb]
\begin{appfig}[graphviz-type / dependency DAG]
% Four ranks at x = 0, 3.6, 7.4, 11.1.  Vertical pitch 1.55, node height 7mm,
% so the minimum gap between two boxes in a rank is 8.5mm and no edge crosses a
% node.  The dashed edge marks the one node with no redundant upstream path.
\node[gvkey] (aw) at (0,0) {Pioneer Award\\five-year term};
\node[gvboxs] (site) at (3.6,1.05) {Site agreement};
\node[gvboxs] (irb) at (3.6,-1.05) {IRB approval and active IND};
\node[gvbox] (bench) at (7.4,1.75) {Bench stop-latency verification};
\node[gvbox] (acc) at (7.4,0.0) {Cohort accrual, 3+3};
\node[gvbox] (mon) at (7.4,-1.75) {Independent safety monitoring};
\node[gvmid] (rp2d) at (11.1,0.95) {RP2D declared};
\node[gvmid] (pub) at (11.1,-0.95) {Public data and code release};
\draw[gvedgeb] (aw) -- (site);
\draw[gvedgeb] (aw) -- (irb);
\draw[gvedge] (site) -- (bench);
\draw[gvedge] (site) -- (acc);
\draw[gvedge] (irb) -- (acc);
\draw[gvedge] (irb) -- (mon);
\draw[gvedge] (bench) -- (rp2d);
\draw[gvedge] (acc) -- (rp2d);
\draw[gvedge] (acc) -- (pub);
\draw[gvedge] (mon) -- (pub);
\node[font=\tiny\itshape,text=protogray,anchor=north,text width=62mm,align=center]
  at (3.6,-1.95) {IND activation is the only node with no redundant upstream
  path, which is why year one funds it first.};
\draw[gvedged] (3.6,-1.85) -- (irb.south);
\end{appfig}
\vspace{-0.7cm}
\figcaption{What the award unblocks, and where the program still has a single
point of failure. Every activity except IND activation has at least two
upstream routes.}
\end{appfloat}
